\chapter{HRBR Signals: a Small Deterministic Reactive Graph}

Part III begins the HRBR runtime. Unlike the VDOM runtime (\texttt{src/*}), HRBR is a "compiled fast path" whose job is to:

\begin{itemize}
  \item track fine-grained dependencies (signals)
  \item schedule updates deterministically (scheduler lanes)
  \item update DOM directly in blocks (later chapters)
\end{itemize}

This chapter focuses on the reactive graph implementation in \texttt{runtime/signals.ts}.

\section{Why signals exist (and why HRBR needs them)}

Up to this point, Relax updates UI via a Virtual DOM pipeline: you re-run \texttt{render()}, diff children, then patch the DOM.
That model is flexible and easy to reason about, but it can do more work than necessary when updates are tiny and frequent.

Signals exist to make a different trade-off:

\begin{quote}
Track \emph{exactly} which computations depend on a value, and when that value changes, re-run only those computations.
\end{quote}

This is the core HRBR idea: the compiler can translate template expressions into tiny computations (effects/memos) that update very specific DOM nodes.
So instead of reconciling a large subtree, a signal write can result in “update this text node”.

Compare the philosophies:

\begin{itemize}
  \item \textbf{VDOM:} change state \textrightarrow rerender \textrightarrow reconcile structure.
  \item \textbf{Signals:} change signal \textrightarrow schedule dependent computations \textrightarrow update exactly where read.
\end{itemize}

For external context (not Relax-specific), see:
\href{https://www.solidjs.com/docs/latest#signals}{SolidJS signals} (fine-grained model) and
\href{https://react.dev/learn/render-and-commit}{React render/commit} (VDOM model).

\section{Public API: the contract we have to satisfy}

The file \texttt{runtime/index.ts} exports these symbols from \texttt{runtime/signals.ts}:

\begin{itemize}
  \item \texttt{createSignal(initial) -> [read, write]}
  \item \texttt{createEffect(fn, options?) -> \{ dispose() \}}
  \item \texttt{createMemo(fn, opts?) -> readMemo}
  \item \texttt{batch(fn) -> result}
  \item \texttt{untrack(fn) -> result}
\end{itemize}

The executable specification lives in:

\begin{itemize}
  \item \texttt{runtime/\_\_tests\_\_/signals.test.ts}
  \item \texttt{runtime/\_\_tests\_\_/signals.determinism.test.ts}
  \item \texttt{runtime/\_\_tests\_\_/signals.fuzz.test.ts}
\end{itemize}

\section{Mental model}

You can read \texttt{runtime/signals.ts} as a deliberately small version of a Solid-like graph:

\begin{itemize}
  \item A \textbf{signal} is a node with a current value and a list of observers.
  \item An \textbf{effect} is a computation that re-runs when its dependencies (signals/memos) change.
  \item A \textbf{memo} is a computation that behaves like a derived signal.
\end{itemize}

Two design constraints keep showing up in the implementation and tests:

\begin{enumerate}
  \item \textbf{Determinism}: observer notification runs in stable creation order.
  \item \textbf{Leak-safety}: disposing a computation must detach it from all sources.
\end{enumerate}

There’s a third, more pragmatic constraint you’ll see in the code:

\begin{itemize}
  \item \textbf{Low ceremony}: everything is implemented with plain objects and a few intrusive linked lists.
\end{itemize}

This is not a ``fully-featured'' reactive system; it’s a compact kernel meant to be embedded into compiled blocks.

\section{The tiny contract (keep this in your head)}

Signals are small, so their contract can be crisp:

\begin{itemize}
  \item \textbf{Read} returns a value and optionally registers a dependency edge.
  \item \textbf{Write} updates the value and schedules observers \emph{in deterministic order}.
  \item \textbf{Effect} runs once initially, then re-runs when its dependencies change.
  \item \textbf{Memo} caches a derived value and notifies dependents when it changes.
  \item \textbf{batch(fn)} coalesces multiple writes into a single downstream flush.
  \item \textbf{untrack(fn)} temporarily disables dependency capture.
\end{itemize}

And the failure modes we must avoid (the tests are built around these):

\begin{itemize}
  \item \textbf{Nondeterminism}: the observer order changes across runs.
  \item \textbf{Leaks}: disposed computations still run after future writes.
  \item \textbf{Re-entrancy recursion}: an effect writes to what it reads and causes stack overflow.
  \item \textbf{Stale dependencies}: conditional reads don’t update subscriptions.
\end{itemize}

\section{Diagram: signals graph (sources \textrightarrow computations \textrightarrow derived values)}
The runtime is easiest to visualize as a graph where reads register edges and writes schedule computations:

\begin{center}
\begin{tikzpicture}[
  node distance=10mm,
  box/.style={draw,rounded corners,fill=RelaxLight,inner sep=5pt,align=center},
  arrow/.style={-Latex,thick}
]
  \node[box] (s) {Signal\\\texttt{s}};
  \node[box, right=22mm of s] (m) {Memo\\\texttt{m = f(s)}};
  \node[box, below=of m] (e) {Effect\\\texttt{e() reads m}};

  \draw[arrow] (s) -- node[above]{read registers observer} (m);
  \draw[arrow] (m) -- node[right]{read registers observer} (e);
  \draw[arrow] (s) to[bend right=24] node[left]{write schedules} (m);
  \draw[arrow] (m) to[bend right=24] node[right]{write schedules} (e);
\end{tikzpicture}
\end{center}

The rest of this chapter explains how Relax keeps this graph deterministic and leak-safe with intrusive linked lists.

\section{Core data structures: intrusive lists for stable order}

In \texttt{runtime/signals.ts} you will see a set of internal types:

\begin{itemize}
  \item \texttt{SignalNode<T>} holds \texttt{value} and an intrusive doubly-linked list of observers: \texttt{observersHead/observersTail}.
  \item \texttt{Computation} represents an effect or memo computation.
  \item \texttt{ObserverLink} is the node stored in the signal's observer list.
  \item \texttt{SourceLink} is a reverse edge stored on the computation so it can detach from sources on re-run/dispose.
\end{itemize}

The key technique: when a computation reads a signal, \texttt{addObserver(node, currentComputation)} inserts the computation into the signal's observer list \emph{sorted by computation id}.

Computation ids come from a global counter:

\begin{verbatim}
let nextId = 0

comp._id = nextId++
\end{verbatim}

Because ids are monotonically increasing, the observer list is stable and deterministic by construction.

\subsection{Why two linked lists (observer list + source list)?}

Each dependency edge is represented twice:

\begin{itemize}
  \item a forward link in the signal’s \texttt{observersHead/observersTail} list (who should be scheduled when the signal changes)
  \item a reverse link in the computation’s \texttt{sourcesHead/sourcesTail} list (what should be detached before re-run/dispose)
\end{itemize}

This dual representation is what makes ``conditional dependencies'' safe: each run starts by detaching previous sources and then re-registering based on the reads that actually occur.

\subsection{Diagram: forward observers vs. reverse sources}

This is the part that makes signals feel like a real runtime (not a toy): every edge is represented twice.

\begin{itemize}
  \item The \textbf{forward edge} lives on the signal: ``who should I notify when I change?''
  \item The \textbf{reverse edge} lives on the computation: ``what did I subscribe to last run, so I can detach?''
\end{itemize}

\begin{center}
\begin{tikzpicture}[
  node distance=8mm,
  box/.style={draw,rounded corners,fill=RelaxLight,inner sep=5pt,align=left},
  arrow/.style={-Latex,thick}
]
  \node[box] (sig) {SignalNode\\\texttt{value}\\\texttt{observersHead/Tail}};
  \node[box, right=26mm of sig] (comp) {Computation\\\texttt{sourcesHead/Tail}\\\texttt{run()} / \texttt{dispose()}};

  \node[box, below=of sig] (ol) {ObserverLink\\(in signal list)\\\texttt{prev/next}};
  \node[box, below=of comp] (sl) {SourceLink\\(in computation list)\\\texttt{prev/next}};

  \draw[arrow] (sig) -- node[above]{forward: notify} (ol);
  \draw[arrow] (comp) -- node[above]{reverse: detach} (sl);
  \draw[arrow] (sl) to[bend left=10] node[below]{points to} (sig);
  \draw[arrow] (sl) to[bend right=10] node[below]{points to} (ol);
\end{tikzpicture}
\end{center}

In \texttt{runtime/signals.ts}, this is the practical outcome:

\begin{itemize}
  \item \texttt{addObserver(node, comp)} inserts an \texttt{ObserverLink} into \texttt{node.observers...}
  \item and it appends a \texttt{SourceLink} into \texttt{comp.sources...} that points back to that observer link
  \item \texttt{unlinkAllSources(comp)} walks \texttt{comp.sources...} and removes each referenced observer link from its signal
\end{itemize}

Senior takeaway:
\begin{quote}
If you don't keep the reverse list, you can't implement conditional dependencies without leaks.
\end{quote}

\section{Dependency capture: \texttt{currentComputation}}

Reactive systems need a way to know "who is currently running".

This file uses a single global pointer:

\begin{verbatim}
let currentComputation: Computation | null
\end{verbatim}

\subsection{Signal reads}

In \texttt{createSignal()}, the reader is:

\begin{itemize}
  \item return \texttt{node.value}
  \item if \texttt{currentComputation != null}, also call \texttt{addObserver(node, currentComputation)}
\end{itemize}

That means dependencies are collected by running the effect/memo function and letting reads "register".

\subsubsection{Mechanics: the global pointer trick}

The \texttt{currentComputation} pointer is the entire core of dependency tracking.
When a computation is running, it sets this global.
Signal reads check it and attach the dependency edge.

Here’s a very small teaching implementation in plain JavaScript (not the real Relax code):

\begin{lstlisting}[language=JavaScript]
let currentComputation = null

function createSignal(initial) {
  const node = { value: initial, observers: [] }

  function read() {
    if (currentComputation) node.observers.push(currentComputation)
    return node.value
  }

  function write(next) {
    const nextValue = typeof next === 'function' ? next(node.value) : next
    if (Object.is(nextValue, node.value)) return
    node.value = nextValue

    for (const comp of node.observers) {
      comp.run()
    }
  }

  return [read, write]
}

function createEffect(fn) {
  const comp = {
    run() {
      const prev = currentComputation
      currentComputation = comp
      try {
        fn()
      } finally {
        currentComputation = prev
      }
    },
  }
  comp.run()
  return { dispose() {} }
}
\end{lstlisting}

Relax’s real implementation adds the hard parts:

\begin{itemize}
  \item leak-safety via a reverse list (so it can detach old edges)
  \item deterministic order via monotonically increasing ids + sorted insertion
  \item correct batching + re-entrancy via a pending set and deterministic flush
\end{itemize}

\paragraph{No dedupe (by design)}

You may notice \texttt{addObserver()} explicitly does not deduplicate observers.
This is safe because every run begins with \texttt{unlinkAllSources(comp)}, which removes all prior edges.
So duplicates can only occur within a single run, and a well-behaved effect/memo typically reads each signal at most once.

\subsection{\texttt{untrack(fn)}}

\texttt{untrack(fn)} temporarily sets \texttt{currentComputation = null} while executing \texttt{fn}.

Spec (\texttt{runtime/\_\_tests\_\_/signals.test.ts}):

\begin{itemize}
  \item \texttt{untrack prevents dependency registration}
\end{itemize}

In other words: you can read a signal without subscribing.

\section{Writes: schedule observers, preserve order}

In \texttt{createSignal()}, the writer:

\begin{enumerate}
  \item computes \texttt{nextValue} (supports "functional set" style)
  \item bails out if \texttt{Object.is(nextValue, node.value)}
  \item updates \texttt{node.value}
  \item walks the observer list and schedules each observer
\end{enumerate}

There’s a subtle but important line in the loop:

\begin{quote}
Capture \texttt{next} first, because running an observer can unsubscribe/resubscribe and mutate the list.
\end{quote}

In code (comments preserved):

\begin{verbatim}
const next = cur.next
scheduleComputation(cur.comp)
cur = next
\end{verbatim}

This ensures we traverse the original linked list in deterministic order.

\subsection*{Diagram: observer notification preserves creation order}
Observer lists are kept in stable creation order (by computation id). A write walks that list and schedules computations without losing determinism:

\begin{center}
\begin{tikzpicture}[
  node distance=7mm,
  box/.style={draw,rounded corners,fill=RelaxLight,inner sep=5pt,align=center},
  arrow/.style={-Latex,thick}
]
  \node[box] (sig) {Signal\\observer list};
  \node[box, below=of sig] (c1) {comp \#1};
  \node[box, below=of c1] (c2) {comp \#2};
  \node[box, below=of c2] (c3) {comp \#3};
  \node[box, right=18mm of c2] (pend) {\texttt{pending}\\(re-entrancy / batch)};

  \draw[arrow] (sig) -- (c1);
  \draw[arrow] (c1) -- (c2);
  \draw[arrow] (c2) -- (c3);
  \draw[arrow] (c2) -- (pend);
\end{tikzpicture}
\end{center}

The key implementation trick is the one shown in the code snippet above: capture \texttt{next} before scheduling so list mutation during computation runs doesn't break traversal.

\paragraph{What gets scheduled?}

When a signal is written, Relax schedules the dependent computations.
Scheduling does not necessarily mean ``run immediately'' (see batching and lanes), but it always preserves the \emph{creation order} contract.

\subsection{A precise code-walk: \texttt{scheduleComputation} and the \texttt{pending} set}

All the tricky correctness properties (re-entrancy, batching, lane scheduling) flow through one choke point:
	\texttt{scheduleComputation(computation)}.

In the real implementation (\texttt{runtime/signals.ts}), there are four gates.
Reading them as a ruleset helps:

\begin{enumerate}
  \item \textbf{Disposed computations are inert.}
  \item \textbf{Running computations don't re-enter.} If a computation schedules itself while running, it is queued in \texttt{pending} and flushed after the current stack unwinds.
  \item \textbf{Lane-aware computations don't run inline.} If a computation has \texttt{\_schedule}, scheduling is delegated to the runtime scheduler.
  \item \textbf{Batching defers work.} If \texttt{isBatching > 0}, add to \texttt{pending} instead of running.
\end{enumerate}

Here's the same logic in a teaching version (simplified JavaScript), preserving the ordering of checks:

\begin{lstlisting}[language=JavaScript]
function scheduleComputation(comp) {
  if (comp._disposed) return

  // Re-entrancy: don't recurse.
  if (comp._running) {
    pending.add(comp)
    return
  }

  // Lane scheduling: hand off to the scheduler.
  if (comp._schedule) {
    comp._schedule(comp)
    return
  }

  // Batching: defer until end of outer batch.
  if (isBatching > 0) {
    pending.add(comp)
    return
  }

  // Default: run inline.
  comp.run()
}
\end{lstlisting}

\subsubsection{Why re-entrancy and batching share the same queue}

At first glance, using one \texttt{pending} set for two distinct situations looks odd.
It's actually a good simplification because both situations mean the same thing:

\begin{quote}
We are not allowed to run this computation right now, but we must not lose the fact that it became dirty.
\end{quote}

The only additional requirement is determinism.
That's why \texttt{flushPendingIfNeeded()} sorts the set by \texttt{\_id} before scheduling.

\subsubsection{Timeline: the re-entrancy convergence test}

Spec: \texttt{re-entrancy: effect that schedules itself via writes does not recurse infinitely}.

The effect is:

\begin{lstlisting}[language=JavaScript]
createEffect(() => {
  const v = s()
  seen.push(v)
  if (v < 3) setS(v + 1)
})
\end{lstlisting}

The important thing is not just that it ends at \texttt{3}, but that it does not recurse.
The run looks like this:

\begin{enumerate}
  \item Start run: comp becomes \texttt{\_running = true}
  \item Reads \texttt{s()} (dependency edge registered)
  \item Calls \texttt{setS(v + 1)} which tries to schedule the same computation
  \item Because it is \texttt{\_running}, scheduling goes to \texttt{pending} (no re-entry)
  \item End run: \texttt{\_running = false}, then we flush \texttt{pending}
\end{enumerate}

This repeats until the signal hits 3.
That is exactly why the expected log is deterministic: \texttt{[0, 1, 2, 3]}.

Reference (conceptual background): Solid's reactivity avoids re-entrancy recursion via queuing and batching as well:
\href{https://www.solidjs.com/docs/latest#execution-order}{Solid: execution order}.

\section{Effects: cleanup, re-run, disposal}

\subsection{Cleanup semantics}

Effects may return an optional cleanup function:

\begin{verbatim}
createEffect(() => {
  ...
  return () => cleanup()
})
\end{verbatim}

The runtime guarantees cleanup runs:

\begin{itemize}
  \item before an effect re-runs
  \item when the effect is disposed
\end{itemize}

Spec: \texttt{effect cleanup runs before re-run and on dispose}.

In code, the cleanup is stored on the computation object as \texttt{comp.cleanup} and executed in a try/finally block both in \texttt{run()} and \texttt{dispose()}.

\subsection{Detaching and resubscribing: \texttt{unlinkAllSources(comp)}}

Before each run, the runtime calls \texttt{unlinkAllSources(comp)}.
That function walks the computation's \texttt{sourcesHead} list and unlinks each corresponding \texttt{ObserverLink} from the signal's observer list.

This gives the common reactive invariant:

\begin{quote}
An effect's dependencies are exactly the signals it read during its last run.
\end{quote}

It also enables conditional dependencies (read signal A in one run, read signal B in another) without leaking observers.

\subsection{Conditional dependencies in slow motion (the reason \texttt{unlinkAllSources} exists)}

Conditional dependencies are the most common place reactive systems silently leak.
The problem looks like this:

\begin{quote}
An effect reads \texttt{a()} in one run, but reads \texttt{b()} in the next run. If you never detach old edges, the effect ends up subscribed to both forever.
\end{quote}

Relax avoids this by making \texttt{unlinkAllSources(comp)} step 0 of every run.
Then the effect/memo re-subscribes only to the signals it \emph{actually reads} this time.

Here's a minimal example to build intuition:

\begin{lstlisting}[language=JavaScript]
const [useA, setUseA] = createSignal(true)
const [a, setA] = createSignal(0)
const [b, setB] = createSignal(0)

createEffect(() => {
  if (useA()) {
    console.log('A path:', a())
  } else {
    console.log('B path:', b())
  }
})
\end{lstlisting}

Now imagine this sequence:

\begin{enumerate}
  \item Initial run: \texttt{useA() == true} so the effect reads \texttt{useA()} and \texttt{a()}.
  \item Flip to \texttt{false}: the effect re-runs, but now it reads \texttt{useA()} and \texttt{b()}.
\end{enumerate}

If the runtime does not detach, a later \texttt{setA(...)} would still schedule the effect even though it no longer reads \texttt{a()}.
That is wasted work and (in large graphs) becomes a performance cliff.

\subsubsection{What the implementation guarantees}

This is the practical contract enforced by the forward+reverse lists:

\begin{itemize}
  \item Before each run, unlink all previously tracked sources.
  \item During the run, every signal/memo read adds new source links.
  \item After the run, the computation is subscribed to exactly the set of reads that occurred.
\end{itemize}

That is why the dual list structure matters: the reverse list is what makes detaching cheap and precise.

Conceptual reference: Solid's explanation of fine-grained dependency tracking is very similar:
\href{https://www.solidjs.com/docs/latest#how-it-works}{Solid: how it works}.

\paragraph{Leak-safety note}

	\texttt{unlinkAllSources()} also deliberately nulls out fields on the link objects (with \texttt{@ts-expect-error} comments).
That’s not required for correctness, but it makes accidental retention less likely in this prototype implementation.

\section{Non-negotiable guarantees (what must remain true)}

If you want to understand signals fast, read the tests first.
They encode the contract in a way that prevents “almost works” implementations.

\subsection{Core semantics: \texttt{runtime/\_\_tests\_\_/signals.test.ts}}

The suite asserts:

\begin{itemize}
  \item dependency tracking (reads register; writes re-run)
  \item \texttt{Object.is} bail-out prevents no-op writes from causing work
  \item \texttt{untrack} disables edge registration
  \item \texttt{batch} is a flush boundary (including nested batching)
  \item cleanup is run before re-run and on dispose
  \item dispose fully detaches
  \item re-entrancy converges without recursive re-entry
  \item non-sync lanes schedule through the runtime scheduler (tests intercept \texttt{setTimeout})
  \item memos behave like derived signals and notify dependents
\end{itemize}

\subsubsection{A senior reading of the suite: what it forces in the implementation}

The most important thing these tests enforce is that signals are not just a pub/sub primitive.
They are a mini runtime with strict invariants:

\begin{itemize}
  \item \textbf{Idempotent writes}: \texttt{Object.is} bail-out means no "useless work" on repeated values.
  \item \textbf{Precise tracking}: \texttt{untrack} demonstrates that dependency capture is opt-in per read, not an ambient global subscription.
  \item \textbf{Deterministic event ordering}: effect order must be stable even under batching.
  \item \textbf{Cleanup correctness}: cleanup must run before the next run, not after.
\end{itemize}

If you've ever debugged a real UI bug caused by nondeterministic ordering, you'll appreciate why determinism is treated as a contract, not an "implementation detail".

\subsection{Determinism + leak-ish checks: \texttt{signals.determinism.test.ts}}

The determinism suite checks a stronger property:
\emph{same seed + same operations produce identical logs across two fresh graphs}.

That forces the runtime to be stable in ways a single unit test won’t catch.
It also includes a practical disposal check: once an effect is disposed, later writes must not produce any new log entries for it.

\subsection{Fuzz/property checks: \texttt{signals.fuzz.test.ts}}

The fuzz test generates nested batches, sets, and disposals.
The primary assertion is determinism (two runs produce identical logs).
There is also a “smoke” property related to disposal: if effects were disposed, later mass-writes shouldn’t explode log volume.

\subsection{Dispose}

Disposing an effect sets \texttt{\_disposed} and detaches from all sources.

Specs:

\begin{itemize}
  \item \texttt{disposed effects fully detach: further writes do not schedule work}
  \item leak-ish checks in \texttt{runtime/\_\_tests\_\_/signals.determinism.test.ts}
\end{itemize}

\section{Re-entrancy: scheduling while running}

In \texttt{scheduleComputation(computation)} there is an explicit re-entrancy rule:

\begin{itemize}
  \item if the computation is \texttt{\_running}, do not recursively call \texttt{run()}
  \item instead, add it to a \texttt{pending} \texttt{Set}
\end{itemize}

Then \texttt{flushPendingIfNeeded()} runs pending computations after the current run unwinds.

\subsection{Two reasons computations are delayed}

The same \texttt{pending} set is used for two distinct delay cases:

\begin{itemize}
  \item \textbf{re-entrancy}: a computation schedules itself while it’s already running
  \item \textbf{batching}: writes should not trigger intermediate runs until the batch boundary
\end{itemize}

Both cases end up in the same deterministic flush code path.

Spec: \texttt{re-entrancy: effect that schedules itself via writes does not recurse infinitely}.

The test constructs an effect that reads a signal and writes it again until a fixed point is reached.
The expected (deterministic) log is:

\begin{verbatim}
[0, 1, 2, 3]
\end{verbatim}

\section{Batching: a flush boundary}

\texttt{batch(fn)} increments \texttt{isBatching} and makes writes collect work into \texttt{pending}.
When the outer-most batch finishes, it flushes the pending set sorted by computation id.

Specs (\texttt{runtime/\_\_tests\_\_/signals.test.ts}):

\begin{itemize}
  \item \texttt{batch groups multiple writes into a single effect re-run}
  \item \texttt{batch() is a flush boundary: nested batches still only produce one re-run after outer finishes}
  \item \texttt{batch() returns the callback result and preserves deterministic effect order after flush}
\end{itemize}

This is the first place you see HRBR's theme: \emph{deterministic scheduling is part of the user-facing contract}.

\paragraph{Implementation detail: deterministic flush}

At batch end, pending computations are flushed by sorting on \texttt{\_id}:

\begin{verbatim}
Array.from(pending).sort((a,b) => a._id - b._id)
\end{verbatim}

This is why the tests can assert stable order like \texttt{['a:2', 'b:2']}.

\section{Memos: derived signals with equality}

\texttt{createMemo(fn, opts?)} creates a computation whose output is stored in a \texttt{SignalNode<T>}.

Important details:

\begin{itemize}
  \item the memo lazily computes on first read (\texttt{if (!hasValue) comp.run()})
  \item it supports an equality function \texttt{opts.equals} (defaults to \texttt{Object.is})
  \item it notifies observers only when the memo value actually changes
\end{itemize}

Spec: \texttt{createMemo computes from signals and updates dependents}.

\paragraph{Memo invalidation path}

When the memo value changes, it walks \texttt{node.observersHead} and schedules each observer computation.
So effects that read a memo are notified exactly like effects that read a base signal.

\section{Lanes: effects can be scheduled via the runtime scheduler}

\texttt{createEffect} accepts options:\\
\texttt{\{ lane?: Lane, budgetMs?: number \}}.

If \texttt{lane === 'sync'} (the default), the initial run executes immediately.

If \texttt{lane !== 'sync'}, \texttt{createEffect} creates a scheduler via \texttt{createScheduler()} (from \texttt{runtime/scheduler.ts}) and schedules runs into that lane.

\subsection{How non-sync lanes are tested}

In \texttt{runtime/\_\_tests\_\_/signals.test.ts}, the tests override \texttt{globalThis.setTimeout} to capture scheduled flush callbacks into an array.
This makes scheduler-driven effects deterministic in unit tests:

\begin{itemize}
  \item createEffect({ lane: 'default' }) does \emph{not} run immediately
  \item the test manually drains the flush queue to trigger the run
\end{itemize}

Specs:

\begin{itemize}
  \item \texttt{createEffect({ lane }) schedules initial run (non-sync)}
  \item \texttt{createEffect forwards budgetMs to scheduler.schedule}
\end{itemize}

The tests intercept \texttt{setTimeout} to make scheduler flushing deterministic in unit tests.

\section{Determinism tests and fuzz tests}

Determinism is explicitly tested at three levels:

\begin{itemize}
  \item Unit test: \texttt{effects run deterministically (creation order)}
  \item Determinism suite: build a random-but-seeded graph and compare logs across fresh instances (\texttt{signals.determinism.test.ts})
  \item Property/fuzz: randomly generate operations (set/batch/dispose) and ensure identical logs for the same seed (\texttt{signals.fuzz.test.ts})
\end{itemize}

These tests are the reason the implementation invests in sorted observer lists and stable flush ordering.

\section{Walking the tests: what they prove}

This chapter is easiest to internalize by reading the tests as executable requirements.

\subsection{Core semantics (\texttt{runtime/\_\_tests\_\_/signals.test.ts})}

\begin{description}
  \item[tracks dependencies and re-runs] Proves that reads register dependencies and writes schedule effects. Also proves the \texttt{Object.is} bail-out on repeated values.
  \item[untrack prevents dependency registration] Proves that \texttt{untrack} temporarily disables dependency capture.
  \item[batch groups writes] Proves that multiple writes inside a batch yield a single re-run.
  \item[nested batch is a flush boundary] Proves \texttt{isBatching} is a counter and only the outermost batch flushes.
  \item[deterministic effect order] Proves flush order is creation order (based on \texttt{\_id}).
  \item[cleanup runs before re-run and on dispose] Proves cleanup semantics for both re-run and disposal.
  \item[disposed effects fully detach] Proves \texttt{unlinkAllSources} works: disposed computations aren’t scheduled by later writes.
  \item[re-entrancy convergence] Proves the pending-set mechanism prevents infinite recursion and produces a deterministic sequence.
  \item[lane scheduling + budgetMs forwarding] Proves non-sync effects are scheduled through the runtime scheduler and that \texttt{budgetMs} is passed through.
  \item[memos] Proves memo recomputation + dependent effect notification.
\end{description}

\subsection{Determinism suite (\texttt{signals.determinism.test.ts})}

This suite tests the strongest property: \emph{same seed + same operation sequence yields the same log across fresh runtime instances}.

It also includes a practical leak-ish check: after disposing effect 0, it mutates signals repeatedly and asserts there are no new log entries starting with \texttt{e0:}.

\subsection{Fuzz/property tests (\texttt{signals.fuzz.test.ts})}

The fuzz suite generates nested batches, writes, and disposals.
The primary assertion is determinism (two independent runs produce identical logs).

It also has a smoke property about disposal effectiveness: if effects were disposed, the final mass-write should not trigger an absurd number of logs (catching obvious detach bugs).

\section{Online resources}

These systems are common in UI runtimes; the references below help connect the dots without assuming the exact same API surface.

\begin{itemize}
  \item SolidJS (conceptual baseline): \href{https://www.solidjs.com/docs/latest#reactivity}{Fine-grained reactivity}
  \item Preact Signals (conceptual): \href{https://preactjs.com/guide/v10/signals/}{Signals and effects}
  \item React (contrast): \href{https://react.dev/learn/render-and-commit}{Render/commit model vs. fine-grained graphs}
  \item A good ``reactivity under the hood'' talk: \href{https://www.youtube.com/watch?v=2H5zYj1fM9U}{SolidJS reactivity deep dive (video)}
  \item Linked list background (for intrusive lists): \href{https://en.wikipedia.org/wiki/Linked_list}{Linked list}
\end{itemize}

\section{Takeaways}

\begin{itemize}
  \item Signals store intrusive observer lists, sorted by computation creation id.
  \item Effects detach and re-subscribe each run, enabling conditional dependencies safely.
  \item Batching and re-entrancy are implemented with a pending set plus deterministic flush ordering.
  \item Lanes integrate with the scheduler so work can be delayed and budgeted.
\end{itemize}

\section{Exercises}

\begin{enumerate}
  \item In \texttt{runtime/signals.ts}, find the exact line in the signal write loop that captures \texttt{next} before scheduling. Explain why this matters when observers unsubscribe/resubscribe while running.
  \item Build a conditional dependency like the \texttt{useA/a/b} example and write down exactly which writes should schedule the effect in each phase.
  \item Using the tests as a guide, explain the difference between \texttt{isBatching} and \texttt{pending}: what does each represent, and why are both needed?
  \item Explain why the \texttt{pending} set is a \texttt{Set} and not an array. What bug would an array allow in batch-heavy code?
\end{enumerate}

\section{Chapter summary}

\begin{itemize}
  \item Signals implement fine-grained dependency tracking via \texttt{currentComputation}: reads register edges.
  \item Determinism is a user-facing guarantee: observers are ordered by creation id, and pending flushes are sorted.
  \item Leak-safety comes from the dual representation of edges (forward observers + reverse sources) and aggressive detaching via \texttt{unlinkAllSources}.
  \item Batching and re-entrancy share the same mechanism: a \texttt{pending} set flushed after the current stack unwinds or after the outer batch ends.
  \item Effects can be lane-scheduled via the HRBR scheduler so the compiler can choose when work runs.
\end{itemize}
